\fontsize{12pt}{14pt} \selectfont
\begin{doublespace}
\chapter{Related Work}

Container security depends on several low-level Linux features that work together to keep applications isolated. These include technologies like namespaces, cgroups, AppArmor, and Seccomp, which help control what a container can see, do, or access. Over the years, researchers and developers have created tools and systems that rely on these features to prevent containers from escaping or causing harm. But these protections are not perfect. Misconfigurations, outdated policies, or changes in the Linux kernel can all lead to serious security issues, such as containers breaking out, gaining extra privileges, or leaking data.

This chapter reviews past work related to container confinement. It covers areas like VM and sandbox escapes, Linux security tools such as AppArmor and cgroups, system call filtering with Seccomp, and monitoring tools built with eBPF. Each section looks at how these tools work, what problems they solve, and where their limitations are. We also explore rule-based security systems and kernel changes that try to improve isolation at a deeper level. Along the way, we explain how Houdini can fill the gaps left by existing approaches.

We finish the chapter with a look at popular tools used for scanning, hardening, and even attacking container setups. These tools show that while many solutions exist, verifying strong isolation still requires hands-on testing.

\section{Virtual Machine Escapes} 

VM escapes occur when an attacker successfully breaks out of a virtualized confined environment, gaining unauthorized access to the host system. Given that hypervisors manage multiple VMs on shared hardware, such escapes can lead to host compromise, cross-VM attacks, and lateral movement within cloud environments. One of the most notable VM escape vulnerabilities is VENOM (CVE-2015-3456) \cite{abdoul2018novel}, which exploited a flaw in the virtual floppy disk controller (FDC) implementation in QEMU, affecting platforms like Xen, KVM, and VirtualBox. Other significant VM escape vulnerabilities include Cloudburst (2009), which exploited VMware's virtual graphics driver; CVE-2018-6981 and CVE-2018-6982, which targeted VMware ESXi, Workstation, and Fusion \cite{vmware2018escape}; and multiple VirtualBox vulnerabilities (2017) enabling guest-to-host code execution \cite{virtualbox2017bugs}. Additionally, side-channel attacks like MDS/ZombieLoad (2019) demonstrated how speculative execution flaws could be leveraged across VMs \cite{mds2019attack}. More recent research has uncovered weaknesses in encrypted virtualization, such as the SEVered attack (2018), which extracted plaintext memory from AMD SEV-protected VMs \cite{severed2018}, and the Heckler attack (2024), which manipulated non-timer interrupts to compromise confidential VMs in AMD SEV-SNP and Intel TDX \cite{heckler2024}.\\

Houdini’s reliance on virtualization acknowledges the potential for VM-level vulnerabilities, but treats the VM boundary as a strong, but not a foolproof line of defense. It uses VMs not as a security guarantee, but as a practical containment layer to isolate test environments from the host. By doing so, Houdini can safely execute untrusted or experimental configurations and observe whether container confinement fails, while minimizing the impact of potential exploitation.

\section{Sandbox Escapes} 

Sandboxing is a critical security mechanism for restricting the execution of untrusted code, commonly used in browsers, mobile apps, and security tools. A sandbox escape occurs when an attacker bypasses these restrictions, executing arbitrary code with higher privileges and potentially leading to full system compromise. Modern web browsers heavily rely on sandboxing to isolate JavaScript execution from the underlying OS. However, several high-profile vulnerabilities have enabled sandbox escapes. For example, CVE-2019-5786 \cite{chromezeroday}, a use-after-free vulnerability in the FileReader API of Google Chrome, allowed attackers to escape Chrome's sandbox and execute code on the host. Other notable sandbox escape vulnerabilities include CVE-2022-26706, which exploited a flaw in macOS's App Sandbox to execute code outside of restricted permissions \cite{macossandboxescape}; CVE-2023-36719, a stack corruption vulnerability in Windows OS that enabled escape from the Chromium sandbox \cite{chromiumsandboxescape}; and CVE-2024-4040, which was actively exploited to break out of CrushFTP’s Virtual File System (VFS) \cite{crushftpsandboxescape}. Additionally, vulnerabilities in Judge0, an online code execution system, exposed systems to full takeover by allowing attackers to escape its sandbox \cite{judge0sandboxescape}. Research has also uncovered critical flaws in vm2, a JavaScript sandbox library, which could be exploited to bypass its security restrictions and gain remote code execution \cite{vm2sandboxescape}.

\section{Process and Resource Isolation}

\subsection{Apparmor}

AppArmor is a Linux security framework based on mandatory access control, similar to SELinux, but it focuses on restricting individual programs \cite{zhu2021lic}. It works by applying predefined security profiles that limit what a program can do \cite{zhu2021lic}. In Docker, AppArmor can be used to load these profiles into containers to restrict their behavior. If no profile is specified, Docker automatically applies a default one \cite{zhu2021lic}.

Mattetti et al. \cite{mattetti2015securing} developed LiCShield as a framework that uses AppArmor as its enforcement backend. It automatically generates custom AppArmor security profiles by analyzing container behavior in a pre-production environment. These profiles are then used to tightly confine both the container’s runtime behavior and its interactions with the host system, ensuring that only authorized actions are permitted during execution.

Sun et al. \cite{sun2018security} extend AppArmor to work independently inside containers by introducing a new concept called a security namespace. Normally, AppArmor is global, it enforces one set of security rules for the whole system, which means containers cannot define or apply their own AppArmor policies. The authors solve this problem by modifying the Linux kernel so that each container can have its own separate AppArmor namespace. This allows a container to load and enforce its own AppArmor profiles for its applications, independent of the host system's policies.

With this change, both the host and the container can apply AppArmor protections at the same time. For example, while the host might enforce a general container confinement policy, the container itself can load a specific profile to restrict what Firefox or MySQL can do \cite{sun2018security}. 

Zhu et al. \cite{zhu2021lic} create Lic-Sec, which is an enhanced AppArmor profile generator that improves Docker container security by combining the strengths of Docker-sec and LiCShield. The reason for creating a similar tool to LiCShield is that LiCShield is incompatible with recent versions of Docker. Lic-Sec automatically creates fine-grained profiles based on container behavior, restricting capabilities, network access, file operations, mounts, and shell execution. By using both Auditd and SystemTap for tracing, Lic-Sec generates stricter, more comprehensive profiles that block privilege escalation and container escape attacks. It adapts to container use cases and enforces these profiles at runtime for stronger and automated protection.

Houdini could have influenced Zhu et al. \cite{zhu2021lic} work by providing an automated framework to evaluate Lic-Sec's effectiveness across diverse environments. For example, instead of manually testing exploits against a fixed kernel version, docker version and set of container images, the authors could have used Houdini to run the same exploit across multiple kernel versions, Docker versions, configurations, and AppArmor profiles to identify gaps in enforcement. Additionally, Houdini’s ability to replay tricks across time and machines would have allowed for more reliable validation of Lic-Sec’s defenses, such as verifying whether blocking /bin/bash consistently prevents privilege escalation across real-world images.



\subsection{Control Groups}

cgroups (short for control groups) are a Linux kernel feature used to manage and limit system resources like CPU, memory, and I/O for groups of processes \cite{chelladhurai2016securing}. They allow administrators or container runtimes like Docker to control how much of a resource each container can use. For example, Docker can apply cgroup settings to restrict a container’s CPU time or memory usage. This helps prevent a single container from consuming too many resources and affecting others on the system \cite{chelladhurai2016securing}. While cgroups are not designed to stop one container from accessing or interfering with another container’s data or processes \cite{chelladhurai2016securing}, they are essential for defending against certain types of Denial-of-Service (DoS) attacks by limiting the damage a container can cause through resource exhaustion \cite{chelladhurai2016securing}.

Chelladhurai et al. \cite{chelladhurai2016securing} identifies DoS attacks as a significant threat in containerized environments like Docker. The authors emphasize that while containers offer process isolation, they lack strong guarantees against resource abuse, which can allow a malicious container to exhaust host resources and impact other containers.

To address this, they propose several practical mitigation strategies using Linux cgroups \cite{chelladhurai2016securing}. They argue that cgroups are essential for defending against DoS attacks by limiting a container’s access to CPU, memory, and I/O resources. However, they also note that cgroups alone do not provide full isolation between containers. Their proposed DoS mitigation technique includes three test cases: setting memory limits directly in the docker run command using -m=512mb, applying default memory limits in the Docker configuration file (/etc/default/docker), and enforcing a system-wide memory limit for containers that do not specify their own \cite{chelladhurai2016securing}. For evaluation, the paper includes sample attack code to simulate DoS scenarios. One example allocates memory indefinitely using malloc, while another consumes CPU resources by endlessly calculating prime numbers. These tests demonstrate how resource limits can effectively prevent containers from overwhelming the host.

The original paper takes a manual approach to testing DoS attacks by running memory and CPU-intensive scripts inside containers with various resource limits. While this method demonstrates the effectiveness of using cgroups to mitigate DoS risks, it is limited in scale and reproducibility. Houdini can automate and greatly expand this type of testing. Instead of running scripts manually, Houdini can encode these attacks as repeatable tricks, such as infinite memory allocation or continuous CPU-bound computations. These tricks can then be executed systematically across a wide range of Docker versions, kernel versions, and container configurations (e.g., with or without memory limits like the -m flag). This is especially important because the implementation and enforcement behavior of cgroups can differ between kernel versions, particularly between cgroups v1 and v2, or in how specific resource limits are handled. By capturing and comparing results, Houdini enables a scalable and repeatable evaluation framework. It can identify which combinations of settings are vulnerable to DoS and which configurations effectively block the attacks, providing a more comprehensive assessment than manual testing.


\subsection{System Call Filtering}

System call filtering is a security technique used to control which system calls a program is allowed to make to the operating system. Since system calls are the primary way that applications interact with the kernel, restricting them helps reduce the attack surface and prevent malicious behavior. On Linux, this is commonly done using Seccomp, which lets developers create filters to allow or block specific syscalls. Traditional filtering with classic BPF (cBPF) is simple and efficient but limited, it only supports static rules and cannot track complex behavior or context. More advanced techniques, such as those using eBPF, improve flexibility by allowing filters to maintain state, inspect arguments, and respond dynamically to syscall patterns. 

This paper used system call filtering as a key part of its security monitoring approach. Gantikow et al. \cite{gantikow2019rule} used two tools named Sysdig and Falco that rely on a kernel module to capture and analyze system calls made by containers. By watching these low-level operations, the system can spot unusual or potentially harmful behavior. Sysdig collects all system calls and lets users filter them by container, process, or syscall type to reduce noise. Falco adds a layer of security by letting users define rules that flag suspicious behavior based on these system calls.

For example, the authors created rules to detect when a container tries to access files outside its expected directories using the open() syscall. They also detected unauthorized programs being run inside containers by monitoring execve calls. Attempts to break out of containers were flagged by watching for the setns syscall, and unexpected network connections were detected by tracking relevant network-related system calls.

While Falco can only detect individual suspicious events, not patterns or sequences of events like those involved in buffer overflows or denial-of-service attacks, the use of system call filtering still proved effective in catching many types of threats with low performance impact.

Ghavamnia et al. \cite{ghavamnia2020confine} present Confine, a system that automatically generates restrictive system call policies for Docker containers. Instead of relying solely on dynamic analysis (which often misses rare or conditional execution paths), Confine uses a hybrid approach: it performs a short dynamic analysis to discover all executables launched during container startup and then uses static analysis on those binaries to extract the complete set of system calls they may invoke. This includes both libc-invoked and directly invoked syscalls.

The goal of this syscall filtering is to allow containers access only to the syscalls they truly need, thereby blocking access to the rest, many of which may contain vulnerabilities or provide attack vectors for privilege escalation or container escape. By creating custom Seccomp BPF profiles that deny all unnecessary syscalls, Confine ensures that even if an attacker gains access to a container, they cannot exploit most of the kernel interface.

The evaluation showed that Confine could disable 145 or more system calls (out of 326) for over half of the tested containers, neutralizing 51 known kernel CVEs. This significantly exceeds Docker’s default Seccomp filter, which only disables 49 syscalls.

Houdini could be used to dynamically test whether Confine’s generated Seccomp profiles are too restrictive or too permissive. For instance, if Houdini can successfully run an attack trick under a Confine policy, it might indicate an over-permissive profile. Conversely, if legitimate workloads fail under Confine’s policy in Houdini’s environment, this suggests over-restriction.

Confine assumes that policies generated statically will remain correct across systems. But Linux syscall implementations and Docker behavior change over time. Houdini could systematically evaluate if the same Confine-generated profile works consistently across kernel versions, and flag new vulnerabilities.

Jiao et al. \cite{jia2023programmable} used system call filtering to improve the security of Linux systems by controlling which system calls a program is allowed to make. Normally, this is done using a tool called Seccomp, which filters system calls based on simple rules written in a limited language called cBPF. However, cBPF is very basic, as it can only check the syscall number and some simple arguments. It cannot remember past behavior, check syscall sequences, or read memory safely. This makes it hard to write strong security rules.

Jiao et al. \cite{jia2023programmable} solve this by creating a new system called Seccomp-eBPF, which uses the more powerful eBPF language instead of cBPF. eBPF allows writing filters that can remember past system calls, count how many times a system call was used, block system calls if they’re too frequent, or only allow them in certain program phases. It can also safely inspect data passed by pointer (like file paths), and prevent known bugs caused by racing system calls.

With these improvements, Seccomp-eBPF can block real-world attacks that cBPF cannot. It also helps reduce the number of system calls allowed in early stages of a program (reducing attack surface) by over 50\% in some cases. The paper shows that this system is both safer and faster than older methods that rely on user-space monitoring, and it works well in environments like containers.


\section{Extended Berkeley Packet Filter}


eBPF (extended Berkeley Packet Filter) is a Linux kernel technology that allows small, sandboxed programs to run safely inside the kernel without modifying kernel source code or loading kernel modules. Originally designed for network packet filtering, eBPF has evolved into a general-purpose tool for tracing, monitoring, and security enforcement. In essence, eBPF lets developers attach custom logic to various kernel events, like system calls, network activity, or file operations, and collect data or enforce rules with minimal performance overhead. Because eBPF programs run in a restricted environment and are verified for safety before execution, they offer a powerful but secure way to extend kernel behavior dynamically.


Lopes et al. \cite{lopes2020container} use eBPF as the core mechanism for tracing system calls made by containerized applications. The goal of the system is to automatically generate custom Seccomp profiles that strictly limit the syscalls a container can perform, thereby reducing the attack surface. To achieve this, the authors use eBPF's powerful kernel-level tracing capabilities to hook into every syscall. Specifically, they implement a BCC (BPF Compiler Collection) program that attaches eBPF probes to all system calls and logs detailed information such as the syscall name, timestamp, process ID, and the namespace from which the syscall originated.

For Lopes et al. \cite{lopes2020container}, the goal is to automatically generate Seccomp profiles by tracing the actual system calls used by a containerized application. While their approach uses eBPF for runtime tracing, Houdini could have provided complementary value by automating and scaling this profiling process across many container configurations, Docker versions, and kernel versions. Houdini’s test matrix model allows systematic exploration of diverse runtime flags (e.g., --privileged, --cap-add, AppArmor profiles) to evaluate whether the generated Seccomp profile remains effective across different settings. Additionally, Houdini could validate whether the reduced syscall profile indeed blocks known exploits, as done in Lopes' Apache test case, but at scale.

This namespace information is crucial because it allows the system to filter out irrelevant syscalls made by the host or other processes, ensuring that only syscalls from the target container are recorded. These logs are then processed to generate a Seccomp profile that whitelists only the syscalls actually used by the application during testing. One of the key advantages of using eBPF is its runtime-agnostic nature. It works not only with Docker but also with Podman or any other container technology, since the tracing happens at the kernel level. Additionally, eBPF's integration into the CI/CD pipeline enables automatic profile generation whenever code changes are committed, making it easy to keep Seccomp profiles up to date with minimal manual effort. As a proof of concept, the authors applied their system to a vulnerable Apache container and demonstrated that, after tracing and generating a Seccomp profile using eBPF, the container was no longer susceptible to a known RCE vulnerability (CVE-2018-11776).

He et al. \cite{he2023cross} found a new type of attack where eBPF can be misused to break container isolation in shared-kernel environments like the cloud. eBPF was designed to safely let programs interact with the Linux kernel for things like monitoring and debugging. But the paper shows that attackers inside a container can use eBPF helper functions, such as ones that read or write other processes’ memory to spy on or interfere with programs running outside the container, including on the host. These attacks can steal sensitive data, crash important services like systemd, or even escape the container by taking over a root-level process on the host.

One key problem is that eBPF tracing programs are not limited by namespaces, so they can see and interact with processes system-wide. The authors tested this in real-world services like Jupyter notebooks and Kubernetes clusters and successfully escaped containers or spread attacks across nodes. Even trusted tools like Cilium and Datadog that use eBPF for monitoring could be at risk if used in containers with the wrong settings.

To fix this, the paper proposes a new system called CapBits. It gives each process fine-grained control over which eBPF features it can use or be targeted by. CapBits is faster and more secure than previous methods and can stop attacks where eBPF is used to hijack sensitive processes. Overall, the paper warns that while eBPF is a powerful tool, its risks in container environments are not well understood and need better protection.

Jiao et al. \cite{jia2023programmable} introduce Seccomp-eBPF, a new way to improve system call filtering on Linux using the extended Berkeley Packet Filter (eBPF). System call filtering is a key technique used to protect the operating system from untrusted applications, such as containers, by controlling which system calls are allowed. The current standard, Seccomp with classic BPF (cBPF), is widely used but has serious limitations. It only allows basic, static checks and cannot store information or respond to complex patterns of behavior, making it hard to implement advanced security policies.

Seccomp-eBPF solves this by allowing developers to write more powerful and flexible filters using eBPF, a newer language already used in many parts of the Linux kernel. With Seccomp-eBPF, filters can track state (such as counting how many times a system call has been used), detect specific sequences of system calls, prevent race-condition attacks by forcing certain system calls to run one at a time, and safely examine data passed from user programs. These features make it possible to block sophisticated attacks that cBPF cannot handle.

The authors implemented this system in the Linux kernel and added support for it in the crun container runtime. They showed that it works well in practice. For example, it reduced the number of allowed system calls during the startup of common server applications by up to 55\%, significantly shrinking the attack surface. Despite its added flexibility, Seccomp-eBPF maintained performance similar to optimized cBPF filters and was much faster than user-space alternatives like Seccomp Notifier. Overall, Seccomp-eBPF offers a safer and more powerful way to secure Linux systems, especially in environments like containers that rely heavily on system call filtering.

Findlay et al. \cite{findlay2021bpfcontain} use eBPF to build a new security system called BPFCONTAIN that improves container security on Linux. Containers often rely on tools like seccomp, AppArmor, and SELinux for isolation, but these tools can be complicated, inconsistent, and not designed specifically for containers. To solve this problem, BPFCONTAIN uses eBPF to monitor and control what containers are allowed to do directly inside the Linux kernel. It checks actions like file access, network use, and system calls, and blocks anything not allowed by the container’s policy. These rules are written in a simple YAML file, making it easy to manage. BPFCONTAIN loads these rules and programs at runtime, without changing the Linux kernel or container tools like Docker. It also adds extra protections to block known container attacks, such as privilege escalation.



\section{Rule-based Systems}

A rule-based system is one that makes decisions or detects events based on a predefined set of rules. These rules are usually written by security experts or administrators and follow simple if-then logic. For example, a rule might state that if a certain system call is made by an unexpected process, the system should flag it as suspicious or block it. This approach allows for clear, predictable behavior because every decision is based on explicitly defined conditions.

In the context of container or system security, rule-based monitoring is often used to detect abnormal behavior. Tools like Falco apply rules to system call traces to identify activities that deviate from what is considered normal or safe. For example, rules can be written to detect unauthorized file access, unusual network activity, or attempts to escalate privileges. These systems are easy to understand and customize, making them widely used in production environments.

However, rule-based systems also have limitations. They can only detect behaviors that have been anticipated and written into the rules. This makes them less effective against new or complex attack patterns that do not match existing rules.

A proposed method by Jian et al. \cite{jian2017defense} implemented as a shell script that runs on the host and follows a sequence of fixed rules to identify suspicious activity. It first collects metadata about all running containers and their processes, including each process’s PID, name, and namespace tag (Nsid). Then, it compares the namespace tag of each container process against the namespace tag of that container’s init process. According to the authors, if a process inside a container has a different namespace tag than its init process, it likely indicates that the process has escaped its namespace—an early sign of a container escape attempt.

This comparison logic is rule-based because it applies a deterministic rule: if Nsid(process) does not equal Nsid(init), then mark the process as suspicious. Further, the script checks whether the suspicious process resembles a shell (based on the process name). If both conditions are met, the system takes action: it kills the process and tries to trace the attacker (e.g., by retrieving the user ID or IP address).

Houdini could have helped improve the work by Jian et al. \cite{jian2017defense} by making the testing more thorough, automated, and easier to repeat. In the paper, the authors propose a way to detect Docker escape attacks by checking if a process inside the container has a different namespace tag than the container’s main process. They test this using only two known attacks: one involving the recvmsg system call and another using the Dirty COW vulnerability.

With Houdini, the authors could have gone much further. Houdini can automatically run many different escape tricks, both known and custom-made across a wide range of Docker versions, Linux kernel versions, and settings like user privileges, Seccomp filters, and namespace configurations. This would help find out whether the detection method works well in more situations or if it fails in certain cases.

Gantikow et al. \cite{gantikow2019rule} propose a rule-based monitoring system that detects security violations in containerized workloads by analyzing system call activity. This is accomplished using two open-source tools: Sysdig and Falco. Sysdig provides a kernel module that captures all system calls made by containers, while Falco builds upon this by applying predefined rules to flag suspicious behavior. These rules are written in a declarative format and can describe conditions such as unauthorized file access, unexpected application launches, or privilege escalation attempts.

In this rule-based model, each rule defines a pattern of behavior that is considered abnormal or malicious. For example, a rule might detect a container process trying to access files outside its designated working directory (indicating a possible information leakage or path traversal attack), or flag execution of unknown binaries inside a container (which might imply a crypto-miner or reverse shell). Rules can be whitelist-based (defining what is allowed) or blacklist-based (defining what is forbidden), and are matched against real-time system call events.

The paper presents various attack scenarios like container breakout, unauthorized networking, or loading kernel modules, and shows how each can be detected with carefully written Falco rules. While the rule-based approach is effective for many types of misuse, the authors acknowledge limitations. Notably, Falco cannot detect frequency-based or multi-event anomalies (e.g., Denial of Service or buffer overflows), as it only supports detection of single-event anomalies.

the paper by Brady et al. \cite{brady2020docker} uses a rule-based system as part of its approach to securing Docker containers. The authors incorporate Anchore Engine into their CI/CD pipeline, which applies a set of predefined security rules, also known as policies, to each Docker image. These rules include checks for known security vulnerabilities and for unsafe Dockerfile configurations, such as whether the image exposes port 22. In their implementation, the authors use Anchore's default rules, but they note that developers can create and apply their own custom rules based on specific security needs. This rule-based system allows the pipeline to automatically detect and block images that do not meet the defined security standards. It helps developers identify and fix issues early in the development process, improving the overall security of containerized applications.

The paper by Zhu et al. \cite{zhu2021lic} describes three tools Docker-sec, LiCShield, and the proposed Lic-Sec all of which rely on generating and enforcing rule-based security profiles using AppArmor. These profiles consist of explicit allow/deny rules that define what operations a containerized application can perform, such as file accesses, network operations, and use of Linux capabilities.

Docker-sec uses dynamic monitoring to learn what capabilities and network accesses a container actually uses, then generates an AppArmor profile with rules that only permit those actions. LiCShield, on the other hand, uses SystemTap to trace kernel-level operations and generates rules for more types of actions, such as file mounts and execution paths. Lic-Sec combines both tools to generate a more comprehensive rule set that includes deny rules for unnecessary shell commands and tighter restrictions on mounts and execution behavior.

In all cases, the tools use observed behavior during a training phase to produce explicit, rule-based security profiles that constrain container behavior at runtime. These rules are enforced by the AppArmor Linux Security Module, making this a clear example of a rule-based system in container security.

The paper by Chen et al. \cite{chen2022seaf} uses a rule-based system as part of its SEAF framework to detect security flaws in Docker container images. SEAF defines a set of detection primitives that act as basic operations for writing flexible detection rules. These primitives include content matching, feature matching, data recording, usage frequency analysis, and file extraction. Each type of vulnerability, such as misconfiguration, unauthorized access, sensitive information leakage, malware, and software vulnerabilities, is detected by combining specific primitives into rules. For example, to detect unauthorized access in Redis containers, SEAF extracts configuration files and uses regular expressions to search for insecure settings like missing passwords. These rule combinations allow SEAF to systematically inspect container layers and identify flaws by applying consistent logic across images. This rule-based approach is modular, meaning new rules can be added easily to address emerging threats without changing the underlying system. As a result, SEAF can adapt to various application scenarios and detect a wide range of security issues using well-defined and reusable rule patterns.

\section{OS Changes}

Luo Yang et al. \cite{luo2016whispers} investigates how containers can still leak information through covert channels, even when isolation mechanisms are in place. It focuses on Docker and shows that the default security features such as namespaces and cgroups are not enough to prevent secret communication between containers. These covert channels take advantage of shared physical resources, allowing containers to bypass the expected security boundaries.

The authors present five high-capacity covert channel techniques. These include CPU load modulation, cache-based covert channels, disk I/O delay signaling, memory bus contention, and file lock contention. Each method allows a sender container to encode data into patterns of resource usage. The receiver container then decodes this data by observing how these patterns affect performance. For example, a sender might create bursts of CPU activity to represent binary data, while the receiver measures CPU latency to interpret the signal. These attacks are a form of side-channel communication, where indirect changes in system behavior are used to leak information.

The paper describes scenarios where a sender container transfers hidden information to a receiver container, despite there being no authorized communication path between them. This demonstrates that containers can communicate by exploiting shared hardware such as CPU caches or memory bandwidth. Even without network access or shared filesystems, the sender can affect system behavior in ways that the receiver can detect and interpret.

In this paper, Sun et al. \cite{sun2018security} introduce the idea of security namespaces, a new feature in the Linux kernel that lets each container have its own security settings using tools like AppArmor and IMA. Normally, these tools are set by the host and apply to all containers. This means containers can’t create or manage their own security rules, which is a problem when you want containers to be more isolated or run different policies. The paper’s solution is to separate these security settings so that each container can have its own, without affecting the host or other containers.

To make this work, the authors made small changes to the Linux kernel, only about 1,000 to 1,500 lines of code. They tested their changes in two ways. First, with IMA, containers were able to set their own policies, like only allowing signed programs to run. The containers correctly blocked unsigned or changed programs, and they couldn’t get around the host’s security rules, which stayed in place. Second, with AppArmor, they tested 20 common programs that already had security profiles. Sixteen of them ran correctly inside containers using both the container’s and the host’s security rules. For the four that didn’t, the system gave warnings ahead of time to avoid crashes.

Even though this new design is useful, it doesn’t work with systems like Houdini. Houdini depends on the way Linux currently handles security, where rules like AppArmor and IMA are shared across the whole system. If Houdini had to work with these new security namespaces, it would need to be redesigned from the ground up. Because of that, this paper’s approach, while creative and effective, isn’t compatible with how Houdini is built today.

The paper by Sun et al. \cite{sun2018security} relate to OS changes because it shows that even though microVMs are meant to be isolated, some of their operations are still handled by the host operating system, creating security risks. The paper finds that parts of the system such as virtiofsd and kernel threads like kvm pit and vhost are not fully isolated. They run with too much access or are not limited by cgroups. To fix this, the authors suggest changing how the OS manages these components by placing kernel threads into the correct cgroups and reducing unnecessary permissions.


\section{Secure Computing}

Lopes et al. \cite{lopes2020container} present a method to improve container security by automating the creation and maintenance of Seccomp profiles. Seccomp is a Linux kernel feature that limits the system calls a containerized application can make, significantly reducing the potential attack surface. However, maintaining these profiles is difficult because it requires a detailed understanding of which syscalls are needed and must be updated whenever the application changes. To address this, the authors develop a system that integrates with CI/CD pipelines to automatically trace the syscalls made during unit and integration testing using eBPF. These syscalls are collected, and a custom Seccomp profile is generated that whitelists only those calls, blocking everything else by default. For applications with incomplete test coverage, the system can use fuzzing to stimulate additional syscall paths. The approach ensures that each container has a tailored Seccomp profile, improving security without manual configuration. The authors show that this method can mitigate real-world vulnerabilities, such as CVE-2018-11776, by blocking the syscalls required for exploitation.

Ghavamnia et al. \cite{ghavamnia2020confine} present a static analysis-based system named Confine that generates restrictive system call policies for Docker containers. Unlike previous solutions that rely on dynamic training with realistic workloads, which can miss rarely executed paths, Confine uses static analysis of container binaries and their dependencies to identify the complete superset of system calls required for correct operation. This allows it to generate Seccomp (Secure Computing Mode) profiles that deny all other system calls, significantly reducing the container’s kernel attack surface.

The authors evaluated Confine on 150 Docker images and found that it could disable over 145 system calls out of 326 for more than half the containers, far more than Docker’s default Seccomp policy. This reduction effectively neutralized 51 known kernel CVEs, demonstrating strong security gains without breaking application functionality. The approach also handles direct syscall invocations, language-specific system call mechanisms such as Go and Java, and dynamic library loading.

Confine directly builds on Seccomp as its enforcement mechanism. It generates custom Seccomp BPF profiles that define allow and deny lists for system calls per container. By statically analyzing what syscalls are truly necessary, Confine tailors these Seccomp profiles to the exact needs of each container, minimizing the set of syscalls exposed to potentially malicious applications and enhancing the host’s protection. This use of Seccomp is central to the paper’s method for reducing kernel attack surfaces and mitigating container escape vulnerabilities.

\section{Practical Tools and Framework Design}

At the time of this writing, to our knowledge, Houdini is the first comprehensive approach for evaluating and comparing \textit{container confinement mechanisms}. Much of the container security literature to date has focused on either vulnerability scanning of container images, building offensive/attack tools for container escapes, or proposing best practices for securing container runtimes. While these papers, tools, and documents do not directly share our research objectives, they still broadly fit into the container security landscape. Additionally, research on VM escapes and sandbox escapes has provided valuable insights into isolation vulnerabilities, many of which relate to container security. Techniques used to break out of hypervisors and sandboxed environments often share similarities with container escape methods, which the importance of strong confinement mechanisms. The remainder of this section does not aim to exhaustively enumerate all container security tools and systems, but rather to highlight the various salient methods, including prior research on times the confinement of VMs and sandboxes have been escaped as it is related to work. Most of the papers in the latter half sections of this chapter were taken from \cite{haq2024sok}

\subsection{Regression Testing.} 

One common practice in continuous integration/continuous deployment (CI/CD) pipelines is the use of automated vulnerability scanning tools to detect known vulnerabilities in container images. Tools such as Trivy \cite{trivy} and Clair \cite{clair} are often integrated into CI pipelines to perform vulnerability assessments during the image build process. These tools automatically scan for security flaws in the container image, ensuring that only secure images are deployed to production. For example, GitLab’s CI/CD pipeline includes built-in features for Container Scanning with Trivy, automatically running vulnerability scans whenever new Docker images are built or updated.

In addition to security scanning and policy enforcement, CI/CD pipelines also use container hardening tools like Docker Bench for Security \cite{docker_bench} and Kube-bench \cite{kube_bench} to ensure containers are securely configured. These tools can be integrated into the pipeline to automatically check container configurations against security best practices, such as ensuring that containers run as non-root users or that unnecessary ports are closed. CircleCI \cite{circleci}, for example, integrates Docker Bench to verify that new container builds follow best practices before they are pushed to production.\\

\subsection{Vulnerability Scanners} 

Many free, open, and closed-source tools exist for identifying the presence of known vulnerabilities in container images. Shu \etal~\cite{shu2017study} developed a vulnerability scanning tool to scan Docker Hub container images at scale. They found that on average, official and community images have concerning amounts of vulnerabilities (180+), and that these vulnerabilities remain unpatched for hundreds of days. Their study utilized the Docker Image Vulnerability Analysis (DIVA) framework, which automates the discovery, download, and analysis of over 300,000 Docker images. The findings revealed that more than 80\% of both official and community images contained at least one high-severity vulnerability, and many images were not updated for extended periods. Moreover, vulnerabilities were often inherited from parent images to child images, further increasing security risks.\\

\subsection{Best Practices} 

Users in search of (often high-level) advice on how to improve the security of their container deployments can consult one of many best practices guides available online. One such guide is the Center for Internet Security (CIS) Docker Benchmark \cite{CIS}. In our experience, guides tend to offer generic advice (e.g., ``\textit{Only allow read access to the root filesystem}''). This is problematic because it oversimplifies container security. Generic advice may seem straightforward, but it doesn't account for the specific context or unique configuration of an individual deployment. Such one-size-fits-all recommendations can lead users to believe that simply following these broad guidelines is enough to secure their systems, potentially leaving critical vulnerabilities unaddressed. In reality, effective container security requires a nuanced, tailored approach that considers the particular needs and threat models of each environment. Houdini fills the gap between generic advice and practical security by enabling one to objectively evaluate the effectiveness of your container confinement measures.\\

\subsection{Offensive Tools} 

The offensive security community has developed a range of specialized tools aimed at identifying and exploiting vulnerabilities within containerized environments. Containers, while providing strong isolation, are not immune to attacks, and these tools focus on testing the robustness of container defenses. CDK\footnote{\url{https://github.com/cdk-team/CDK}} and DEEPCE\footnote{\url{https://github.com/stealthcopter/deepce}} are two widely used, open-source container penetration testing toolkits that leverage a collection of known exploits to gain persistence, escape the container, and gather sensitive information about the container environment. These tools aim to bypass isolation mechanisms like namespaces, cgroups, and SELinux, often using a variety of aggressive techniques to identify weaknesses in container security. CDK specifically uses any available method to circumvent security measures, providing a broad toolkit to test the limits of a container's defense. Meanwhile, DEEPCE, developed by stealthcopter, focuses on privilege escalation and container enumeration, attempting to exploit specific weaknesses to achieve container escapes, thereby compromising the host system. In contrast, Houdini takes a more systematic and structured approach to container security testing. Rather than focusing solely on exploiting vulnerabilities, Houdini allows for the direct comparison of defensive techniques by clearly defining the testing environment, outlining the exploit steps, and documenting the expected outputs.\\


\end{doublespace}